\begin{abstract}
Autonomous knowledge systems that operate over extended timescales require a
reliable, self-regulating infrastructure layer.  We present the \Core{}, a
reusable, event-driven substrate that provides four key capabilities for
self-monitoring: a high-throughput event bus with persistence, adaptive
thresholds calibrated to system complexity, a chaos detector with circuit
breaker and predictive warnings, and a thermodynamic cost model that ties
computational effort to carbon impact.  Each component is implemented in
Python, fully testable, and designed to be wired together via the event bus.
We evaluate the core on a series of synthetic workloads that stress its
self-regulation mechanisms: event throughput under back-pressure, threshold
adaptation to drifting complexity, early detection of oscillatory instability,
and energy-aware rejection of inefficient computations.  On an ASUS Strix
GL704GW laptop (Intel Core i7‑8750H, NVIDIA RTX 2070), the event bus
sustains 121\,k events/s for small messages with a per‑event overhead of
$6.9\,\mu$s.  A predictive warning mechanism detects divergent behaviour 78
time steps earlier than a conventional threshold detector, and GPU‑based
training consumes 24$\times$ less energy than the CPU equivalent.  The
\Core{} forms the execution substrate of the larger \Apeiron{} Framework and
is released as open-source software.
\end{abstract}